\documentclass[12pt, a4paper]{article}

% ─── Packages ────────────────────────────────────────────────────────────────
\usepackage[utf8]{inputenc}
\usepackage[T1]{fontenc}
\usepackage[margin=2.5cm, top=3cm, bottom=3cm, headheight=28pt]{geometry}
\usepackage[table]{xcolor}
\usepackage{fancyhdr}
\usepackage{titlesec}
\usepackage{booktabs}
\usepackage{longtable}
\usepackage{array}
\usepackage{enumitem}
\usepackage{hyperref}
\usepackage{tabularx}
\usepackage{multirow}
\usepackage{makecell}
\usepackage{parskip}

% ─── Colors ──────────────────────────────────────────────────────────────────
\definecolor{primary}{RGB}{27, 58, 107}
\definecolor{secondary}{RGB}{46, 95, 172}
\definecolor{rowgray}{RGB}{244, 246, 250}
\definecolor{musthave}{RGB}{198, 228, 179}
\definecolor{shouldhave}{RGB}{255, 242, 204}
\definecolor{nicetohave}{RGB}{222, 222, 222}
\definecolor{risklow}{RGB}{198, 228, 179}
\definecolor{riskmed}{RGB}{255, 242, 204}

% ─── Hyperref ─────────────────────────────────────────────────────────────────
\hypersetup{
    colorlinks=true,
    linkcolor=secondary,
    urlcolor=secondary,
    pdftitle={SRS - Traccia Travel Organizer Offline},
    pdfauthor={Gruppo XX},
}

% ─── Header / Footer ─────────────────────────────────────────────────────────
\pagestyle{fancy}
\fancyhf{}
\fancyhead[L]{\small\color{secondary}SRS -- Traccia Travel Organizer Offline}
\fancyhead[R]{\small\color{secondary}Gruppo XX}
\fancyfoot[C]{\small\thepage}
\renewcommand{\headrulewidth}{0.4pt}
\renewcommand{\headrule}{\hbox to\headwidth{\color{secondary}\leaders\hrule height \headrulewidth\hfill}}

% ─── Section styles ──────────────────────────────────────────────────────────
\titleformat{\section}
  {\Large\bfseries\color{primary}}{\thesection}{1em}{}
  [\vspace{2pt}{\color{secondary}\hrule height 1.2pt}\vspace{4pt}]

\titleformat{\subsection}
  {\large\bfseries\color{secondary}}{\thesubsection}{1em}{}

\titleformat{\subsubsection}
  {\normalsize\bfseries\color{primary}}{\thesubsubsection}{1em}{}

% ─── Table helpers ────────────────────────────────────────────────────────────
\newcolumntype{L}[1]{>{\raggedright\arraybackslash}p{#1}}
\newcolumntype{C}[1]{>{\centering\arraybackslash}p{#1}}
\renewcommand{\arraystretch}{1.3}

% ─── Req-block macro (replicates reference SRS style) ────────────────────────
% Usage: \reqblock{Codice}{Categoria}{Priorità}{Rischio}{UC correlati}
\newcommand{\reqblock}[5]{%
  \begin{tabular}{@{}ll@{}}
    \toprule
    \textbf{Codice:}              & #1 \\
    \textbf{Categoria:}           & #2 \\
    \textbf{Priorità:}            & #3 \\
    \textbf{Rischio tecnico:}     & #4 \\
    \textbf{Casi d'uso correlati:}& #5 \\
    \bottomrule
  \end{tabular}
}

% Priority macros
\newcommand{\MUST}{\colorbox{musthave}{\small\textbf{MUST HAVE}}}
\newcommand{\SHOULD}{\colorbox{shouldhave}{\small\textbf{SHOULD HAVE}}}
\newcommand{\NICE}{\colorbox{nicetohave}{\small\textbf{NICE TO HAVE}}}
\newcommand{\RISKL}{\colorbox{risklow}{\small Basso}}
\newcommand{\RISKM}{\colorbox{riskmed}{\small Medio}}

% ─── UC-block macro ───────────────────────────────────────────────────────────
% Provides consistent use-case header
\newcommand{\ucblock}[3]{%
  \noindent
  \begin{tabular}{@{}ll@{}}
    \textbf{Attore:}         & #1 \\
    \textbf{Precondizioni:}  & #2 \\
    \textbf{Postcondizioni:} & #3 \\
  \end{tabular}%
  \medskip
}

% =============================================================================
\begin{document}
% =============================================================================

% ─── Title Page ──────────────────────────────────────────────────────────────
\begin{titlepage}
  \centering
  \vspace*{1.8cm}
  {\color{primary}\rule{\linewidth}{1.5pt}}\\[0.5cm]
  {\Huge\bfseries\color{primary}
    Documento di Specifica\\[0.3cm]
    dei Requisiti Software}\\[0.5cm]
  {\color{primary}\rule{\linewidth}{1.5pt}}\\[1cm]
  {\LARGE\bfseries\color{secondary}Traccia --- Travel Organizer Offline App}\\[0.8cm]
  {\large Corso di Ingegneria del Software --- Gruppo XX}\\[0.3cm]
  {\large Dicembre 2025}\\[2.5cm]

  \begin{tabular}{ll}
    \toprule
    \textbf{Documento} & Specifica dei Requisiti Software (SRS) \\
    \textbf{Versione}  & 1.0 \\
    \textbf{Stato}     & Definitivo \\
    \textbf{Data}      & Dicembre 2025 \\
    \textbf{Gruppo}    & XX \\
    \bottomrule
  \end{tabular}
  \vfill
\end{titlepage}

% ─── Indice ──────────────────────────────────────────────────────────────────
\newpage
\tableofcontents
\newpage

% =============================================================================
\section{Descrizione dell'Applicazione}
% =============================================================================

\subsection{Scopo dell'applicazione}

Traccia è un'applicazione mobile sviluppata in React Native che permette di organizzare
viaggi in modo completo, mantenendo tutte le informazioni disponibili anche senza connessione
internet. L'idea nasce da un'esigenza concreta: quando si è in viaggio la connessione
non è sempre garantita, ma la necessità di consultare l'itinerario, aggiornare le spese
o spuntare elementi dalla valigia non aspetta. L'app gestisce l'intero ciclo di vita di
un viaggio, dalla pianificazione iniziale fino al consuntivo delle spese al rientro.

\subsection{Funzionalità principali}

\begin{itemize}[leftmargin=1.5em, itemsep=2pt]
  \item Gestione completa dei viaggi con stati di avanzamento e budget
  \item Suddivisione del viaggio in tappe o giornate con itinerario ordinato
  \item Pianificazione di attività associate a tappe, date e luoghi specifici
  \item Checklist personalizzabili per oggetti, documenti e cose da fare
  \item Gestione delle spese con distinzione tra previsto ed effettivo
  \item Sezione statistiche con confronto budget e distribuzione delle spese
  \item Timeline giornaliera dell'itinerario (feature avanzata)
  \item Duplicazione di una tappa con tutte le attività associate (feature avanzata)
  \item Persistenza locale completa: nessun account, nessun server, nessuna rete
\end{itemize}

\subsection{Attori coinvolti}

\textbf{Utente:} la persona che pianifica e gestisce i propri viaggi attraverso l'app.
Non è previsto alcun sistema di autenticazione né di differenziazione di ruoli.
Tutti i dati appartengono al dispositivo su cui l'app è installata.

% =============================================================================
\section{Requisiti Funzionali}
% =============================================================================

\subsection{Business Flow}

\subsubsection{BF-1: Gestione completa dei viaggi}

\begin{tabular}{@{}ll@{}}
  \toprule
  \textbf{Priorità:}        & \MUST \\
  \textbf{Rischio tecnico:} & \RISKL \\
  \bottomrule
\end{tabular}

\medskip
\noindent
L'utente deve poter gestire l'intero ciclo di vita dei propri viaggi, dall'inserimento
iniziale alla cancellazione definitiva, passando per la modifica e la consultazione.
Ogni viaggio aggrega tutte le informazioni ad esso collegate: tappe, attività, checklist
e spese. La lista dei viaggi deve essere consultabile e filtrabile in modo rapido.

\subsubsection{BF-2: Gestione delle tappe e dell'itinerario}

\begin{tabular}{@{}ll@{}}
  \toprule
  \textbf{Priorità:}        & \MUST \\
  \textbf{Rischio tecnico:} & \RISKL \\
  \bottomrule
\end{tabular}

\medskip
\noindent
L'utente deve poter suddividere un viaggio in tappe o giornate, organizzate in ordine
cronologico. Ogni tappa rappresenta l'unità strutturale attorno a cui si aggregano le
attività. L'itinerario deve essere consultabile in modo lineare e permettere la navigazione
diretta verso le attività di ciascuna tappa.

\subsubsection{BF-3: Gestione delle attività}

\begin{tabular}{@{}ll@{}}
  \toprule
  \textbf{Priorità:}        & \MUST \\
  \textbf{Rischio tecnico:} & \RISKM \\
  \bottomrule
\end{tabular}

\medskip
\noindent
L'utente deve poter pianificare attività specifiche -- visite, pasti, spostamenti,
prenotazioni, eventi e altro -- associate a un viaggio e, opzionalmente, a una tappa.
Il sistema deve permettere di tracciare lo stato di ciascuna attività e di filtrarle
per categoria o per giornata.

\subsubsection{BF-4: Gestione di checklist e informazioni utili}

\begin{tabular}{@{}ll@{}}
  \toprule
  \textbf{Priorità:}        & \MUST \\
  \textbf{Rischio tecnico:} & \RISKL \\
  \bottomrule
\end{tabular}

\medskip
\noindent
L'utente deve poter creare e gestire checklist legate al viaggio per tenere traccia di
oggetti da portare, documenti da preparare o attività da completare prima della partenza.
Gli elementi delle checklist possono essere marcati come completati e il progresso è
visualizzato in modo sintetico. Oltre alle checklist, il sistema supporta note libere
e promemoria collegati al viaggio.

\subsubsection{BF-5: Gestione delle spese}

\begin{tabular}{@{}ll@{}}
  \toprule
  \textbf{Priorità:}        & \MUST \\
  \textbf{Rischio tecnico:} & \RISKL \\
  \bottomrule
\end{tabular}

\medskip
\noindent
L'utente deve poter registrare le spese del viaggio distinguendo tra importi previsti
e importi effettivamente sostenuti. Il sistema calcola automaticamente i totali e li
confronta con il budget definito in fase di creazione del viaggio. Le spese sono
categorizzabili e filtrabili.

\subsubsection{BF-6: Ricerca, filtri e statistiche}

\begin{tabular}{@{}ll@{}}
  \toprule
  \textbf{Priorità:}        & \SHOULD \\
  \textbf{Rischio tecnico:} & \RISKL \\
  \bottomrule
\end{tabular}

\medskip
\noindent
Il sistema deve offrire funzionalità di ricerca e filtraggio sui principali dati gestiti,
nonché una schermata di riepilogo con statistiche aggregate. Le statistiche comprendono
conteggi per stato di viaggio, avanzamento delle attività, distribuzione delle spese per
categoria e confronto tra budget previsto e spesa effettiva.

% =============================================================================
\subsection{Funzionalità Individuali}

\subsubsection{IF-1.1: Inserimento nuovo viaggio}

\medskip
\reqblock{IF-1.1}{Funzionalità}{\MUST}{\RISKL}{UC-1}
\medskip

\noindent
Il sistema deve permettere la creazione di un nuovo viaggio specificando i seguenti dati:
titolo, destinazione principale, data di inizio, data di fine, descrizione o note generali,
stato iniziale, budget previsto, lista dei partecipanti e informazioni utili. Il sistema
deve validare che la data di fine non sia antecedente alla data di inizio e che il titolo
non sia vuoto. Lo stato iniziale predefinito è \textit{Pianificato}.

\subsubsection{IF-1.2: Modifica dati viaggio}

\medskip
\reqblock{IF-1.2}{Funzionalità}{\MUST}{\RISKL}{UC-2}
\medskip

\noindent
Il sistema deve permettere la modifica di tutti i campi di un viaggio esistente. Le stesse
regole di validazione applicate in fase di inserimento si applicano anche alla modifica.
Il sistema deve aggiornare i dati e mostrare un messaggio di conferma al completamento.

\subsubsection{IF-1.3: Cancellazione viaggio}

\medskip
\reqblock{IF-1.3}{Funzionalità}{\MUST}{\RISKL}{UC-3}
\medskip

\noindent
Il sistema deve permettere la cancellazione di un viaggio previa richiesta di conferma
esplicita. La cancellazione è a cascata: vengono eliminati automaticamente tutte le tappe,
le attività, le checklist e le spese associate al viaggio eliminato.

\subsubsection{IF-1.4: Visualizzazione lista viaggi}

\medskip
\reqblock{IF-1.4}{Funzionalità}{\MUST}{\RISKL}{UC-4}
\medskip

\noindent
Il sistema deve visualizzare l'elenco di tutti i viaggi ordinato per data di inizio in
ordine decrescente. Per ciascun viaggio vengono mostrati titolo, destinazione, date di
inizio e fine, stato e budget previsto. Nel caso in cui non siano presenti viaggi salvati,
il sistema mostra un messaggio informativo con un invito alla creazione del primo viaggio.

\subsubsection{IF-1.5: Ricerca e filtro viaggi}

\medskip
\reqblock{IF-1.5}{Funzionalità}{\MUST}{\RISKL}{UC-5}
\medskip

\noindent
Il sistema deve permettere la ricerca testuale per destinazione (ricerca parziale,
case-insensitive) e il filtraggio per stato del viaggio. I risultati si aggiornano
in tempo reale durante la digitazione. Se la ricerca non produce risultati, il sistema
mostra il messaggio \textit{``Nessun viaggio trovato''}.

\subsubsection{IF-2.1: Inserimento nuova tappa}

\medskip
\reqblock{IF-2.1}{Funzionalità}{\MUST}{\RISKL}{UC-6}
\medskip

\noindent
Il sistema deve permettere l'aggiunta di una tappa a un viaggio specificando: titolo,
data, località, descrizione e note. Il titolo è obbligatorio. L'ordine delle tappe
nell'itinerario è determinato automaticamente dalla data e può essere modificato
manualmente dall'utente.

\subsubsection{IF-2.2: Modifica tappa}

\medskip
\reqblock{IF-2.2}{Funzionalità}{\MUST}{\RISKL}{UC-7}
\medskip

\noindent
Il sistema deve permettere la modifica di tutti i dati di una tappa esistente.
Le stesse regole di validazione applicate in fase di inserimento rimangono valide.

\subsubsection{IF-2.3: Cancellazione tappa}

\medskip
\reqblock{IF-2.3}{Funzionalità}{\MUST}{\RISKL}{UC-8}
\medskip

\noindent
Il sistema deve permettere la cancellazione di una tappa previa richiesta di conferma.
La cancellazione è a cascata: vengono eliminate automaticamente tutte le attività
associate alla tappa rimossa.

\subsubsection{IF-2.4: Visualizzazione tappe del viaggio}

\medskip
\reqblock{IF-2.4}{Funzionalità}{\MUST}{\RISKL}{UC-9}
\medskip

\noindent
Il sistema deve visualizzare le tappe di un viaggio ordinate per data. Per ciascuna
tappa vengono mostrati titolo, data, località e numero di attività associate.
Selezionando una tappa l'utente accede al dettaglio con le attività collegate.

\subsubsection{IF-3.1: Inserimento nuova attività}

\medskip
\reqblock{IF-3.1}{Funzionalità}{\MUST}{\RISKM}{UC-10}
\medskip

\noindent
Il sistema deve permettere l'aggiunta di un'attività specificando: titolo (obbligatorio),
descrizione, data e fascia oraria, luogo, categoria, costo previsto, stato e note.
L'attività deve essere collegata a un viaggio e, opzionalmente, a una tappa. Le categorie
previste sono: Visita, Pasto, Escursione, Spostamento, Prenotazione, Evento e Altro.
Lo stato predefinito è \textit{Da svolgere}. Il costo previsto, se specificato, deve
essere un valore numerico maggiore o uguale a zero.

\subsubsection{IF-3.2: Modifica attività}

\medskip
\reqblock{IF-3.2}{Funzionalità}{\MUST}{\RISKL}{UC-11}
\medskip

\noindent
Il sistema deve permettere la modifica di tutti i campi di un'attività esistente,
incluso lo stato (Da svolgere, Completata, Annullata).

\subsubsection{IF-3.3: Cancellazione attività}

\medskip
\reqblock{IF-3.3}{Funzionalità}{\MUST}{\RISKL}{UC-12}
\medskip

\noindent
Il sistema deve permettere la cancellazione di un'attività previa richiesta di conferma.
La cancellazione di un'attività non comporta l'eliminazione a cascata di altri dati.

\subsubsection{IF-3.4: Visualizzazione e filtro attività}

\medskip
\reqblock{IF-3.4}{Funzionalità}{\MUST}{\RISKL}{UC-13}
\medskip

\noindent
Il sistema deve permettere di visualizzare le attività di un viaggio o di una tappa
specifica, con la possibilità di filtrarle per categoria e per stato. Le attività
vengono visualizzate ordinate per data e ora.

\subsubsection{IF-4.1: Gestione delle checklist}

\medskip
\reqblock{IF-4.1}{Funzionalità}{\MUST}{\RISKL}{UC-14}
\medskip

\noindent
Il sistema deve permettere di creare una o più checklist per ogni viaggio, di aggiungere,
modificare ed eliminare elementi e di marcarli come completati o da completare. Il sistema
mostra il progresso di ogni checklist come rapporto tra elementi completati e totali
(ad esempio 4/10). La cancellazione di una checklist elimina anche tutti i suoi elementi.

\subsubsection{IF-4.2: Gestione delle informazioni utili}

\medskip
\reqblock{IF-4.2}{Funzionalità}{\SHOULD}{\RISKL}{UC-15}
\medskip

\noindent
Il sistema deve permettere di aggiungere, modificare ed eliminare note libere, promemoria
e riferimenti a prenotazioni associati al viaggio, come numeri di telefono utili, indirizzi
di alloggi o orari di voli e treni.

\subsubsection{IF-5.1: Inserimento nuova spesa}

\medskip
\reqblock{IF-5.1}{Funzionalità}{\MUST}{\RISKL}{UC-16}
\medskip

\noindent
Il sistema deve permettere la registrazione di una spesa specificando: titolo (obbligatorio),
importo (obbligatorio, maggiore di zero), categoria, data, tipo (Prevista o Effettiva),
metodo di pagamento e note. La spesa deve essere collegata al viaggio e, opzionalmente,
a una tappa o a un'attività. Le categorie previste sono: Alloggio, Trasporto, Cibo,
Attività, Shopping e Altro.

\subsubsection{IF-5.2: Modifica e cancellazione spesa}

\medskip
\reqblock{IF-5.2}{Funzionalità}{\MUST}{\RISKL}{UC-17, UC-18}
\medskip

\noindent
Il sistema deve permettere la modifica di tutti i campi di una spesa esistente e la
sua cancellazione previa richiesta di conferma.

\subsubsection{IF-5.3: Visualizzazione e filtro spese}

\medskip
\reqblock{IF-5.3}{Funzionalità}{\MUST}{\RISKL}{UC-19}
\medskip

\noindent
Il sistema deve visualizzare le spese di un viaggio mostrando il totale delle spese
previste e il totale delle spese effettive, con indicazione dello scostamento rispetto
al budget. L'utente può filtrare le spese per categoria e per tipo (Prevista/Effettiva).

\subsubsection{IF-6.1: Riepilogo e statistiche}

\medskip
\reqblock{IF-6.1}{Funzionalità}{\SHOULD}{\RISKL}{UC-20}
\medskip

\noindent
Il sistema deve fornire una schermata di riepilogo con le seguenti informazioni: numero
totale di viaggi per stato, attività pianificate e completate per il viaggio selezionato,
spesa totale prevista e spesa totale effettiva con confronto rispetto al budget, distribuzione
delle spese per categoria, progresso complessivo delle checklist. Le informazioni numeriche
sono affiancate da rappresentazioni grafiche a barre o a torta.

% =============================================================================
\subsection{Requisiti relativi ai Dati}

\subsubsection{DF-1.1: Struttura dati Viaggio}

\begin{tabular}{@{}ll@{}}
  \toprule
  \textbf{Codice:}          & DF-1.1 \\
  \textbf{Categoria:}       & Dati \\
  \textbf{Priorità:}        & \MUST \\
  \textbf{Rischio tecnico:} & \RISKL \\
  \bottomrule
\end{tabular}

\medskip
\noindent
Il sistema deve gestire per ogni viaggio i seguenti dati:

\begin{itemize}[leftmargin=1.5em, itemsep=1pt]
  \item id (stringa, univoco, generato automaticamente)
  \item Titolo (stringa, obbligatorio, max 100 caratteri)
  \item Destinazione (stringa, obbligatorio, max 100 caratteri)
  \item Data di inizio (data, obbligatorio)
  \item Data di fine (data, obbligatorio, \(\geq\) data di inizio)
  \item Descrizione (stringa, opzionale, max 500 caratteri)
  \item Stato (enumerazione: Pianificato, In corso, Completato, Archiviato)
  \item Budget previsto (decimale, opzionale, \(\geq 0\))
  \item Partecipanti (lista di stringhe, opzionale)
  \item Informazioni utili (testo libero, opzionale)
\end{itemize}

\subsubsection{DF-1.2: Struttura dati Tappa}

\begin{tabular}{@{}ll@{}}
  \toprule
  \textbf{Codice:}          & DF-1.2 \\
  \textbf{Categoria:}       & Dati \\
  \textbf{Priorità:}        & \MUST \\
  \textbf{Rischio tecnico:} & \RISKL \\
  \bottomrule
\end{tabular}

\medskip
\noindent
Il sistema deve gestire per ogni tappa i seguenti dati:

\begin{itemize}[leftmargin=1.5em, itemsep=1pt]
  \item id (stringa, univoco, generato automaticamente)
  \item tripId (riferimento al viaggio, obbligatorio)
  \item Titolo (stringa, obbligatorio, max 100 caratteri)
  \item Data (data, obbligatorio)
  \item Località (stringa, opzionale, max 100 caratteri)
  \item Descrizione (stringa, opzionale, max 500 caratteri)
  \item Ordine (intero, calcolato automaticamente per data, modificabile manualmente)
  \item Note (stringa, opzionale)
\end{itemize}

\subsubsection{DF-1.3: Struttura dati Attività}

\begin{tabular}{@{}ll@{}}
  \toprule
  \textbf{Codice:}          & DF-1.3 \\
  \textbf{Categoria:}       & Dati \\
  \textbf{Priorità:}        & \MUST \\
  \textbf{Rischio tecnico:} & \RISKL \\
  \bottomrule
\end{tabular}

\medskip
\noindent
Il sistema deve gestire per ogni attività i seguenti dati:

\begin{itemize}[leftmargin=1.5em, itemsep=1pt]
  \item id (stringa, univoco, generato automaticamente)
  \item tripId (riferimento al viaggio, obbligatorio)
  \item stageId (riferimento alla tappa, opzionale)
  \item Titolo (stringa, obbligatorio, max 100 caratteri)
  \item Descrizione (stringa, opzionale, max 500 caratteri)
  \item Data e ora (datetime, opzionale)
  \item Luogo (stringa, opzionale, max 100 caratteri)
  \item Categoria (enumerazione: Visita, Pasto, Escursione, Spostamento, Prenotazione, Evento, Altro)
  \item Costo previsto (decimale, opzionale, \(\geq 0\))
  \item Stato (enumerazione: Da svolgere, Completata, Annullata)
  \item Note (stringa, opzionale)
\end{itemize}

\subsubsection{DF-1.4: Struttura dati Spesa}

\begin{tabular}{@{}ll@{}}
  \toprule
  \textbf{Codice:}          & DF-1.4 \\
  \textbf{Categoria:}       & Dati \\
  \textbf{Priorità:}        & \MUST \\
  \textbf{Rischio tecnico:} & \RISKL \\
  \bottomrule
\end{tabular}

\medskip
\noindent
Il sistema deve gestire per ogni spesa i seguenti dati:

\begin{itemize}[leftmargin=1.5em, itemsep=1pt]
  \item id (stringa, univoco, generato automaticamente)
  \item tripId (riferimento al viaggio, obbligatorio)
  \item stageId (riferimento alla tappa, opzionale)
  \item activityId (riferimento all'attività, opzionale)
  \item Titolo (stringa, obbligatorio, max 100 caratteri)
  \item Importo (decimale, obbligatorio, \(> 0\))
  \item Categoria (enumerazione: Alloggio, Trasporto, Cibo, Attività, Shopping, Altro)
  \item Data (data, obbligatorio)
  \item Tipo (enumerazione: Prevista, Effettiva)
  \item Metodo di pagamento (stringa, opzionale)
  \item Note (stringa, opzionale)
\end{itemize}

\subsubsection{DF-1.5: Struttura dati Checklist e ChecklistItem}

\begin{tabular}{@{}ll@{}}
  \toprule
  \textbf{Codice:}          & DF-1.5 \\
  \textbf{Categoria:}       & Dati \\
  \textbf{Priorità:}        & \MUST \\
  \textbf{Rischio tecnico:} & \RISKL \\
  \bottomrule
\end{tabular}

\medskip
\noindent
Il sistema deve gestire i seguenti dati per le checklist:

\begin{itemize}[leftmargin=1.5em, itemsep=1pt]
  \item id (stringa, univoco)
  \item tripId (riferimento al viaggio, obbligatorio)
  \item Titolo (stringa, obbligatorio, max 100 caratteri)
\end{itemize}

\noindent
E i seguenti per ciascun elemento della checklist:

\begin{itemize}[leftmargin=1.5em, itemsep=1pt]
  \item id (stringa, univoco)
  \item checklistId (riferimento alla checklist, obbligatorio)
  \item Testo (stringa, obbligatorio, max 200 caratteri)
  \item Completato (booleano, valore predefinito: false)
\end{itemize}

\subsubsection{DF-2.1: Persistenza dei dati}

\begin{tabular}{@{}ll@{}}
  \toprule
  \textbf{Codice:}          & DF-2.1 \\
  \textbf{Categoria:}       & Dati \\
  \textbf{Priorità:}        & \MUST \\
  \textbf{Rischio tecnico:} & \RISKM \\
  \bottomrule
\end{tabular}

\medskip
\noindent
Tutti i dati gestiti dall'app -- viaggi, tappe, attività, checklist e spese -- devono
essere salvati in modo persistente sul dispositivo tramite AsyncStorage (per React Native)
o Hive/Isar (per Flutter). I dati vengono caricati all'avvio dell'app e salvati
automaticamente dopo ogni operazione di scrittura. L'app deve funzionare correttamente
in assenza di connessione di rete in ogni sua parte.

% =============================================================================
\subsection{Requisiti dell'Interfaccia Utente}

\subsubsection{UI-1.1: Navigazione principale}

\begin{tabular}{@{}ll@{}}
  \toprule
  \textbf{Codice:}    & UI-1.1 \\
  \textbf{Categoria:} & Interfaccia \\
  \textbf{Priorità:}  & \MUST \\
  \bottomrule
\end{tabular}

\medskip
\noindent
L'app deve presentare una navigazione coerente e intuitiva tra le schermate principali,
con accesso alla lista dei viaggi dalla schermata iniziale e navigazione gerarchica
verso tappe, attività, checklist e spese di ciascun viaggio. La schermata delle
statistiche deve essere raggiungibile da qualsiasi punto dell'app.

\subsubsection{UI-1.2: Form di inserimento e modifica}

\begin{tabular}{@{}ll@{}}
  \toprule
  \textbf{Codice:}    & UI-1.2 \\
  \textbf{Categoria:} & Interfaccia \\
  \textbf{Priorità:}  & \MUST \\
  \bottomrule
\end{tabular}

\medskip
\noindent
I form devono presentare campi chiaramente etichettati con input appropriati per tipo
(tasto testo, selettore data, campo numerico, menu a tendina). La validazione avviene
in tempo reale con messaggi di errore posizionati sotto al campo interessato.

\subsubsection{UI-1.3: Tabelle e liste di visualizzazione dati}

\begin{tabular}{@{}ll@{}}
  \toprule
  \textbf{Codice:}    & UI-1.3 \\
  \textbf{Categoria:} & Interfaccia \\
  \textbf{Priorità:}  & \MUST \\
  \bottomrule
\end{tabular}

\medskip
\noindent
Le liste devono mostrare le informazioni principali di ciascun elemento con opzioni
di azione rapida accessibili tramite pressione prolungata o menu contestuale. Ogni
schermata deve gestire lo stato vuoto con un messaggio informativo.

\subsubsection{UI-1.4: Ricerca e filtri}

\begin{tabular}{@{}ll@{}}
  \toprule
  \textbf{Codice:}    & UI-1.4 \\
  \textbf{Categoria:} & Interfaccia \\
  \textbf{Priorità:}  & \MUST \\
  \bottomrule
\end{tabular}

\medskip
\noindent
Le schermate di lista che supportano la ricerca devono presentare una barra di ricerca
nella parte superiore con aggiornamento in tempo reale. I filtri per categoria o stato
sono presentati come chip o menu a tendina sopra la lista.

\subsubsection{UI-1.5: Indicatori visivi di stato}

\begin{tabular}{@{}ll@{}}
  \toprule
  \textbf{Codice:}    & UI-1.5 \\
  \textbf{Categoria:} & Interfaccia \\
  \textbf{Priorità:}  & \SHOULD \\
  \bottomrule
\end{tabular}

\medskip
\noindent
Lo stato delle attività (Da svolgere, Completata, Annullata), delle spese (Prevista,
Effettiva) e dei viaggi deve essere evidenziato tramite colori e icone coerenti in
tutta l'app. Lo scostamento negativo rispetto al budget deve essere segnalato in rosso.

\subsubsection{UI-1.6: Messaggi di conferma e di errore}

\begin{tabular}{@{}ll@{}}
  \toprule
  \textbf{Codice:}    & UI-1.6 \\
  \textbf{Categoria:} & Interfaccia \\
  \textbf{Priorità:}  & \MUST \\
  \bottomrule
\end{tabular}

\medskip
\noindent
Il sistema deve mostrare un dialogo di conferma prima di ogni operazione distruttiva
e messaggi di errore chiari e leggibili in caso di validazione fallita o operazione
non consentita.

\subsubsection{UI-2.1: Lingua dell'interfaccia}

\begin{tabular}{@{}ll@{}}
  \toprule
  \textbf{Codice:}    & UI-2.1 \\
  \textbf{Categoria:} & Interfaccia \\
  \textbf{Priorità:}  & \SHOULD \\
  \bottomrule
\end{tabular}

\medskip
\noindent
L'interfaccia utente è interamente in lingua italiana.

% =============================================================================
\section{Requisiti Non Funzionali}
% =============================================================================

\subsubsection{FC-1: Affidabilità e gestione degli errori}

\begin{tabular}{@{}ll@{}}
  \toprule
  \textbf{Codice:}          & FC-1 \\
  \textbf{Categoria:}       & Affidabilità \\
  \textbf{Priorità:}        & \MUST \\
  \textbf{Rischio tecnico:} & \RISKM \\
  \bottomrule
\end{tabular}

\medskip
\noindent
Il sistema deve:
\begin{itemize}[leftmargin=1.5em, itemsep=1pt]
  \item Gestire tutti gli errori derivanti da operazioni asincrone senza causare crash dell'app
  \item Mostrare messaggi di errore chiari e comprensibili in caso di operazione fallita
  \item Non perdere dati già salvati in caso di errore su un'operazione successiva
\end{itemize}

\subsubsection{FC-2: Integrità dei dati}

\begin{tabular}{@{}ll@{}}
  \toprule
  \textbf{Codice:}          & FC-2 \\
  \textbf{Categoria:}       & Integrità \\
  \textbf{Priorità:}        & \MUST \\
  \textbf{Rischio tecnico:} & \RISKM \\
  \bottomrule
\end{tabular}

\medskip
\noindent
Il sistema deve garantire che:
\begin{itemize}[leftmargin=1.5em, itemsep=1pt]
  \item La cancellazione di un viaggio rimuova a cascata tutte le entità collegate
  \item I riferimenti tra entità (stageId in un'attività, tripId in una spesa) siano
    sempre validi oppure esplicitamente nulli
  \item Non sia possibile salvare dati che violino i vincoli definiti nei requisiti DF
\end{itemize}

\subsubsection{FC-3: Validazione dei dati}

\begin{tabular}{@{}ll@{}}
  \toprule
  \textbf{Codice:}          & FC-3 \\
  \textbf{Categoria:}       & Qualità dei dati \\
  \textbf{Priorità:}        & \MUST \\
  \textbf{Rischio tecnico:} & \RISKL \\
  \bottomrule
\end{tabular}

\medskip
\noindent
Tutti i dati inseriti devono essere validati prima del salvataggio:
\begin{itemize}[leftmargin=1.5em, itemsep=1pt]
  \item Campi obbligatori: verifica non vuoti
  \item Range di valori: importo \(> 0\), costo attività \(\geq 0\), data fine \(\geq\) data inizio
  \item Lunghezze massime: rispettare i limiti specificati nei requisiti DF
  \item Formati: date nel formato corretto
\end{itemize}

\subsubsection{FC-4: Prestazioni}

\begin{tabular}{@{}ll@{}}
  \toprule
  \textbf{Codice:}          & FC-4 \\
  \textbf{Categoria:}       & Prestazioni \\
  \textbf{Priorità:}        & \SHOULD \\
  \textbf{Rischio tecnico:} & \RISKL \\
  \bottomrule
\end{tabular}

\medskip
\noindent
Il sistema deve garantire:
\begin{itemize}[leftmargin=1.5em, itemsep=1pt]
  \item Operazioni CRUD: completamento in meno di 1 secondo
  \item Rendering di liste fino a 100 elementi: meno di 2 secondi
  \item Caricamento iniziale dell'app: meno di 3 secondi
\end{itemize}

\subsubsection{FC-5: Portabilità}

\begin{tabular}{@{}ll@{}}
  \toprule
  \textbf{Codice:}          & FC-5 \\
  \textbf{Categoria:}       & Portabilità \\
  \textbf{Priorità:}        & \SHOULD \\
  \textbf{Rischio tecnico:} & \RISKL \\
  \bottomrule
\end{tabular}

\medskip
\noindent
L'app deve essere eseguibile su Android (API level 21 e superiori) e su iOS 13 e
superiori. In caso di utilizzo di React Native, deve poter essere eseguita sia su
emulatore che su dispositivo fisico mediante Expo o React Native CLI.

\subsubsection{FC-6: Manutenibilità}

\begin{tabular}{@{}ll@{}}
  \toprule
  \textbf{Codice:}          & FC-6 \\
  \textbf{Categoria:}       & Manutenibilità \\
  \textbf{Priorità:}        & \SHOULD \\
  \textbf{Rischio tecnico:} & \RISKL \\
  \bottomrule
\end{tabular}

\medskip
\noindent
Il codice deve seguire i principi di Separation of Concerns e Single Responsibility.
I componenti riutilizzabili devono essere estratti in moduli condivisi e non duplicati.
La struttura delle cartelle e le scelte architetturali devono essere documentate nel file
README della repository.

\subsubsection{FC-7: Funzionamento offline}

\begin{tabular}{@{}ll@{}}
  \toprule
  \textbf{Codice:}          & FC-7 \\
  \textbf{Categoria:}       & Disponibilità \\
  \textbf{Priorità:}        & \MUST \\
  \textbf{Rischio tecnico:} & \RISKL \\
  \bottomrule
\end{tabular}

\medskip
\noindent
Tutte le funzionalità devono essere pienamente accessibili in assenza di connessione
di rete. Non è previsto alcun backend remoto né servizio esterno di qualsiasi tipo.

% =============================================================================
\section{Scelte Tecnologiche e di Implementazione}
% =============================================================================

\subsection{Framework e linguaggio}

Il progetto è sviluppato in \textbf{React Native} con \textbf{TypeScript}. La scelta
di React Native rispetto a Flutter è motivata da tre fattori principali. Prima di tutto
il gruppo ha già familiarità con JavaScript ed il passaggio a TypeScript comporta
una curva di apprendimento più contenuta rispetto a Dart. In secondo luogo l'ecosistema
npm offre un numero elevato di librerie mature per le funzionalità richieste dalla
traccia, in particolare per la persistenza locale e la visualizzazione di grafici.
Infine Expo semplifica notevolmente il setup iniziale e il testing su dispositivi fisici
senza richiedere la configurazione di ambienti nativi.

TypeScript è preferito al JavaScript puro perché la tipizzazione statica riduce
significativamente la classe di errori a runtime, soprattutto nella gestione dei dati
persistenti dove un campo mancante o di tipo errato può causare comportamenti difficili
da debuggare.

\subsection{Navigazione}

Per la navigazione tra schermate si utilizza \textbf{React Navigation 6} (libreria
\texttt{@react-navigation/native}). La struttura di navigazione è organizzata su due livelli:

\begin{itemize}[leftmargin=1.5em, itemsep=2pt]
  \item \textbf{Tab Navigator} in fondo allo schermo con tre voci principali:
    Viaggi, Statistiche e Impostazioni
  \item \textbf{Stack Navigator} per ciascun tab, che gestisce la navigazione
    gerarchica verso i dettagli (viaggio \(\to\) tappa \(\to\) attività, ecc.)
\end{itemize}

I dati vengono passati tra le schermate tramite i parametri della route di React
Navigation per oggetti semplici, oppure recuperati direttamente dallo store globale
per le schermate che necessitano di dati aggregati.

\subsection{Gestione dello stato}

La gestione dello stato è affidata a \textbf{Zustand}, una libreria leggera che non
richiede la quantità di codice boilerplate di Redux e si integra bene con TypeScript.
Lo stato è suddiviso in cinque store separati -- uno per ciascuna entità principale
(viaggi, tappe, attività, checklist, spese) -- ognuno con i propri selettori e azioni.
Questa suddivisione riduce i re-render inutili e facilita il testing isolato dei singoli
moduli di stato.

\subsection{Persistenza locale}

La persistenza è gestita tramite \textbf{AsyncStorage} (\texttt{@react-native-async-storage/async-storage}).
I dati sono serializzati in JSON e salvati con chiavi gerarchiche del tipo
\texttt{trips}, \texttt{stages:\{tripId\}}, \texttt{activities:\{tripId\}} e così via.
Il layer di persistenza è incapsulato in un servizio dedicato (\texttt{StorageService})
che espone metodi asincroni di lettura e scrittura. Gli store Zustand si occupano di
chiamare il servizio ad ogni modifica, rendendo la persistenza trasparente per i
componenti UI.

Si è valutata l'alternativa di SQLite tramite \texttt{expo-sqlite} per query più
complesse, ma la struttura dei dati dell'app non richiede join relazionali complessi
e AsyncStorage risulta sufficiente per i volumi attesi.

\subsection{Librerie principali}

\begin{itemize}[leftmargin=1.5em, itemsep=3pt]
  \item \texttt{@react-navigation/native} + \texttt{@react-navigation/stack} +
    \texttt{@react-navigation/bottom-tabs} --- navigazione tra schermate
  \item \texttt{zustand} --- gestione dello stato globale
  \item \texttt{@react-native-async-storage/async-storage} --- persistenza locale
  \item \texttt{react-native-gifted-charts} --- grafici per la schermata statistiche
    (grafici a torta e a barre con animazioni)
  \item \texttt{@react-native-community/datetimepicker} --- selettore nativo di
    date e ore nei form
  \item \texttt{react-native-uuid} --- generazione di identificatori univoci per le entità
  \item \texttt{react-native-vector-icons} --- icone per la navigazione e gli stati
\end{itemize}

\subsection{Organizzazione del codice}

Il progetto segue una struttura a feature-based folders: ciascun dominio (viaggi,
tappe, attività, checklist, spese, statistiche) ha la propria cartella con schermate,
componenti specifici e logica di business. Il codice condiviso (componenti riutilizzabili,
tipi TypeScript, servizi, utilità) è collocato in cartelle dedicate al di fuori delle
feature. Questa organizzazione facilita la suddivisione del lavoro tra i membri del
gruppo e riduce i conflitti di merge.

\begin{verbatim}
src/
  features/
    trips/        (schermate e componenti dei viaggi)
    stages/       (schermate e componenti delle tappe)
    activities/   (schermate e componenti delle attività)
    checklists/   (schermate e componenti delle checklist)
    expenses/     (schermate e componenti delle spese)
    stats/        (schermata statistiche)
  shared/
    components/   (componenti UI riutilizzabili)
    hooks/        (custom hook condivisi)
    store/        (store Zustand)
    services/     (StorageService e altri servizi)
    types/        (definizioni TypeScript delle entità)
    utils/        (funzioni di utilità)
  navigation/     (configurazione React Navigation)
\end{verbatim}

% =============================================================================
\section{Casi d'Uso}
% =============================================================================

\subsection{Diagrammi dei Casi d'Uso}

I casi d'uso sono organizzati in tre diagrammi tematici per mantenere la leggibilità.

\bigskip
\begin{center}
\fbox{\begin{minipage}{0.88\textwidth}
  \centering
  \vspace{2cm}
  \textit{[Inserire qui il Diagramma UC --- Gruppo 1: Gestione Viaggi]}\\[0.4cm]
  \small
  Attore: Utente\\
  UC-1 Crea Viaggio \quad UC-2 Modifica Viaggio \quad UC-3 Elimina Viaggio\\
  UC-4 Visualizza Lista Viaggi \quad UC-5 Cerca e Filtra Viaggi
  \vspace{2cm}
\end{minipage}}

\medskip
{\small Figura 1: Diagramma dei casi d'uso -- Gestione Viaggi}
\end{center}

\bigskip
\begin{center}
\fbox{\begin{minipage}{0.88\textwidth}
  \centering
  \vspace{2cm}
  \textit{[Inserire qui il Diagramma UC --- Gruppo 2: Tappe, Attività e Checklist]}\\[0.4cm]
  \small
  Attore: Utente\\
  UC-6 Aggiungi Tappa \quad UC-7 Modifica Tappa \quad UC-8 Elimina Tappa \quad UC-9 Visualizza Tappe\\
  UC-10 Aggiungi Attività \quad UC-11 Modifica Attività \quad UC-12 Elimina Attività \quad UC-13 Filtra Attività\\
  UC-14 Gestisci Checklist \quad UC-15 Gestisci Informazioni Utili\\
  UC-21 Timeline Giornaliera \textsf{«include»} UC-13 \quad
  UC-22 Duplica Tappa \textsf{«include»} UC-6
  \vspace{2cm}
\end{minipage}}

\medskip
{\small Figura 2: Diagramma dei casi d'uso -- Tappe, Attività e Checklist}
\end{center}

\bigskip
\begin{center}
\fbox{\begin{minipage}{0.88\textwidth}
  \centering
  \vspace{2cm}
  \textit{[Inserire qui il Diagramma UC --- Gruppo 3: Spese e Statistiche]}\\[0.4cm]
  \small
  Attore: Utente\\
  UC-16 Aggiungi Spesa \quad UC-17 Modifica Spesa \quad UC-18 Elimina Spesa\\
  UC-19 Visualizza e Filtra Spese \quad UC-20 Visualizza Statistiche e Riepilogo
  \vspace{2cm}
\end{minipage}}

\medskip
{\small Figura 3: Diagramma dei casi d'uso -- Spese e Statistiche}
\end{center}

% =============================================================================
\subsection{Descrizione dettagliata dei Casi d'Uso}

% ─── UC-1 ───────────────────────────────────────────────────────────────────
\subsubsection{UC-1: Crea Viaggio}

\ucblock{Utente}{L'utente si trova nella schermata Lista Viaggi}{Il nuovo viaggio è salvato e appare nell'elenco}

\noindent\textbf{Flusso normale:}
\begin{enumerate}[leftmargin=1.5em, itemsep=1pt]
  \item L'utente seleziona ``Nuovo viaggio''
  \item Il sistema mostra il form di inserimento con i campi del viaggio
  \item L'utente compila i campi obbligatori (titolo, destinazione, date)
  \item L'utente conferma l'inserimento
  \item Il sistema valida i dati inseriti
  \item Il sistema salva il viaggio nello storage locale
  \item Il sistema torna alla lista e mostra il viaggio appena creato
\end{enumerate}

\noindent\textbf{Flussi alternativi:}
\begin{itemize}[leftmargin=1.5em, itemsep=1pt]
  \item 5a. Dati non validi (campo obbligatorio vuoto, data fine antecedente a data inizio):
    il sistema evidenzia i campi con errore e mostra un messaggio descrittivo;
    l'esecuzione riprende dal passo 3.
\end{itemize}

% ─── UC-2 ───────────────────────────────────────────────────────────────────
\subsubsection{UC-2: Modifica Viaggio}

\ucblock{Utente}{Esiste almeno un viaggio; l'utente si trova nel dettaglio viaggio}{I dati del viaggio sono aggiornati}

\noindent\textbf{Flusso normale:}
\begin{enumerate}[leftmargin=1.5em, itemsep=1pt]
  \item L'utente seleziona l'opzione ``Modifica'' dal menu del viaggio
  \item Il sistema mostra il form precompilato con i dati attuali
  \item L'utente modifica uno o più campi
  \item L'utente conferma le modifiche
  \item Il sistema valida i nuovi dati
  \item Il sistema aggiorna il viaggio nello storage e mostra conferma
\end{enumerate}

\noindent\textbf{Flussi alternativi:}
\begin{itemize}[leftmargin=1.5em, itemsep=1pt]
  \item 5a. Dati non validi: il sistema segnala gli errori, l'esecuzione riprende dal passo 3.
\end{itemize}

% ─── UC-3 ───────────────────────────────────────────────────────────────────
\subsubsection{UC-3: Elimina Viaggio}

\ucblock{Utente}{Esiste almeno un viaggio}{Il viaggio e tutti i dati associati sono eliminati}

\noindent\textbf{Flusso normale:}
\begin{enumerate}[leftmargin=1.5em, itemsep=1pt]
  \item L'utente seleziona l'opzione ``Elimina'' dal menu del viaggio
  \item Il sistema mostra un dialogo di conferma con avviso sull'eliminazione a cascata
  \item L'utente conferma
  \item Il sistema rimuove il viaggio con tutte le tappe, attività, checklist e spese collegate
  \item Il sistema torna alla lista viaggi
\end{enumerate}

\noindent\textbf{Flussi alternativi:}
\begin{itemize}[leftmargin=1.5em, itemsep=1pt]
  \item 3a. L'utente annulla: il sistema chiude il dialogo senza apportare modifiche.
\end{itemize}

% ─── UC-4 ───────────────────────────────────────────────────────────────────
\subsubsection{UC-4: Visualizza Lista Viaggi}

\ucblock{Utente}{L'utente avvia l'app o naviga alla schermata principale}{La lista dei viaggi è visualizzata}

\noindent\textbf{Flusso normale:}
\begin{enumerate}[leftmargin=1.5em, itemsep=1pt]
  \item Il sistema carica i dati dallo storage locale
  \item Il sistema ordina i viaggi per data di inizio in ordine decrescente
  \item Il sistema visualizza titolo, destinazione, date e stato per ciascun viaggio
  \item Se non esistono viaggi il sistema mostra il messaggio ``Nessun viaggio salvato.
    Crea il tuo primo viaggio!''
\end{enumerate}

% ─── UC-5 ───────────────────────────────────────────────────────────────────
\subsubsection{UC-5: Cerca e Filtra Viaggi}

\ucblock{Utente}{L'utente si trova nella schermata Lista Viaggi}{I viaggi corrispondenti ai criteri sono visualizzati}

\noindent\textbf{Flusso normale:}
\begin{enumerate}[leftmargin=1.5em, itemsep=1pt]
  \item L'utente inserisce testo nella barra di ricerca oppure seleziona un filtro di stato
  \item Il sistema aggiorna la lista in tempo reale mostrando i risultati corrispondenti
  \item Se nessun viaggio soddisfa i criteri il sistema mostra ``Nessun viaggio trovato''
\end{enumerate}

% ─── UC-6 ───────────────────────────────────────────────────────────────────
\subsubsection{UC-6: Aggiungi Tappa}

\ucblock{Utente}{L'utente si trova nella sezione Tappe di un viaggio}{La nuova tappa è salvata e appare nell'itinerario}

\noindent\textbf{Flusso normale:}
\begin{enumerate}[leftmargin=1.5em, itemsep=1pt]
  \item L'utente seleziona ``Aggiungi tappa''
  \item Il sistema mostra il form con i campi: titolo, data, località, descrizione, note
  \item L'utente compila i campi obbligatori e conferma
  \item Il sistema valida i dati e salva la tappa
  \item Il sistema aggiorna l'itinerario ordinando le tappe per data
\end{enumerate}

\noindent\textbf{Flussi alternativi:}
\begin{itemize}[leftmargin=1.5em, itemsep=1pt]
  \item 4a. Titolo vuoto: il sistema segnala l'errore, l'esecuzione riprende dal passo 2.
\end{itemize}

% ─── UC-7 / UC-8 ─────────────────────────────────────────────────────────────
\subsubsection{UC-7: Modifica Tappa}

\ucblock{Utente}{Esiste almeno una tappa}{I dati della tappa sono aggiornati}

\noindent Flusso analogo a UC-2 (Modifica Viaggio), applicato alla tappa.

\subsubsection{UC-8: Elimina Tappa}

\ucblock{Utente}{Esiste almeno una tappa}{La tappa e tutte le attività associate sono eliminate}

\noindent Flusso analogo a UC-3 (Elimina Viaggio). La cancellazione a cascata riguarda
le attività collegate alla tappa rimossa.

% ─── UC-9 ────────────────────────────────────────────────────────────────────
\subsubsection{UC-9: Visualizza Tappe del Viaggio}

\ucblock{Utente}{L'utente si trova nel dettaglio di un viaggio}{Le tappe del viaggio sono visualizzate in ordine cronologico}

\noindent\textbf{Flusso normale:}
\begin{enumerate}[leftmargin=1.5em, itemsep=1pt]
  \item Il sistema recupera le tappe del viaggio dallo storage
  \item Il sistema le ordina per data
  \item Il sistema visualizza titolo, data, località e numero di attività associate per ciascuna
  \item Selezionando una tappa l'utente accede al dettaglio con le attività collegate
\end{enumerate}

% ─── UC-10 ───────────────────────────────────────────────────────────────────
\subsubsection{UC-10: Aggiungi Attività}

\ucblock{Utente}{L'utente si trova nella sezione Attività di un viaggio o di una tappa}{La nuova attività è salvata}

\noindent\textbf{Flusso normale:}
\begin{enumerate}[leftmargin=1.5em, itemsep=1pt]
  \item L'utente seleziona ``Aggiungi attività''
  \item Il sistema mostra il form con i campi: titolo, descrizione, data e ora,
    luogo, categoria, costo previsto, stato e note
  \item L'utente compila i campi e conferma
  \item Il sistema valida i dati (titolo obbligatorio, costo \(\geq 0\) se specificato)
  \item Il sistema salva l'attività e aggiorna la lista
\end{enumerate}

\noindent\textbf{Flussi alternativi:}
\begin{itemize}[leftmargin=1.5em, itemsep=1pt]
  \item 4a. Titolo vuoto o costo negativo: il sistema segnala l'errore,
    l'esecuzione riprende dal passo 2.
\end{itemize}

% ─── UC-11 / UC-12 ───────────────────────────────────────────────────────────
\subsubsection{UC-11: Modifica Attività}

\ucblock{Utente}{Esiste almeno un'attività}{I dati dell'attività sono aggiornati}

\noindent Flusso analogo a UC-2. Comprende anche la modifica dello stato
(Da svolgere, Completata, Annullata).

\subsubsection{UC-12: Elimina Attività}

\ucblock{Utente}{Esiste almeno un'attività}{L'attività è eliminata}

\noindent Flusso analogo a UC-3. La cancellazione di un'attività non comporta
eliminazione a cascata di altri dati.

% ─── UC-13 ───────────────────────────────────────────────────────────────────
\subsubsection{UC-13: Visualizza e Filtra Attività}

\ucblock{Utente}{L'utente si trova nella sezione Attività di un viaggio o di una tappa}{Le attività filtrate sono visualizzate}

\noindent\textbf{Flusso normale:}
\begin{enumerate}[leftmargin=1.5em, itemsep=1pt]
  \item Il sistema mostra tutte le attività del viaggio o della tappa, ordinate per data e ora
  \item L'utente seleziona un filtro per categoria e/o per stato
  \item Il sistema aggiorna la lista in tempo reale
\end{enumerate}

% ─── UC-14 ───────────────────────────────────────────────────────────────────
\subsubsection{UC-14: Gestisci Checklist}

\ucblock{Utente}{L'utente si trova nella sezione Checklist di un viaggio}{La checklist è aggiornata e il progresso è salvato}

\noindent\textbf{Flusso normale:}
\begin{enumerate}[leftmargin=1.5em, itemsep=1pt]
  \item Il sistema visualizza le checklist del viaggio con il relativo progresso
  \item L'utente può creare una nuova checklist inserendo il titolo
  \item L'utente aggiunge elementi alla checklist inserendo il testo
  \item L'utente può modificare il testo di un elemento, eliminarlo o marcarlo come completato
  \item Ogni modifica viene salvata immediatamente nello storage locale
\end{enumerate}

% ─── UC-15 ───────────────────────────────────────────────────────────────────
\subsubsection{UC-15: Gestisci Informazioni Utili}

\ucblock{Utente}{L'utente si trova nella sezione Info del viaggio}{Le note sono salvate}

\noindent L'utente può aggiungere, modificare ed eliminare note libere, promemoria o
riferimenti utili collegati al viaggio. Ogni modifica viene salvata immediatamente.

% ─── UC-16 ───────────────────────────────────────────────────────────────────
\subsubsection{UC-16: Aggiungi Spesa}

\ucblock{Utente}{L'utente si trova nella sezione Spese di un viaggio}{La nuova spesa è salvata e i totali sono aggiornati}

\noindent\textbf{Flusso normale:}
\begin{enumerate}[leftmargin=1.5em, itemsep=1pt]
  \item L'utente seleziona ``Aggiungi spesa''
  \item Il sistema mostra il form con i campi: titolo, importo, categoria, data,
    tipo (Prevista/Effettiva), metodo di pagamento e note
  \item L'utente compila i campi obbligatori e conferma
  \item Il sistema valida i dati (titolo obbligatorio, importo \(> 0\)) e salva
  \item Il sistema aggiorna il totale e la lista delle spese
\end{enumerate}

\noindent\textbf{Flussi alternativi:}
\begin{itemize}[leftmargin=1.5em, itemsep=1pt]
  \item 4a. Importo non valido (\(\leq 0\) o non numerico): il sistema segnala l'errore,
    l'esecuzione riprende dal passo 2.
\end{itemize}

% ─── UC-17 / UC-18 ───────────────────────────────────────────────────────────
\subsubsection{UC-17: Modifica Spesa}

\ucblock{Utente}{Esiste almeno una spesa}{I dati della spesa sono aggiornati}

\noindent Flusso analogo a UC-2 (Modifica Viaggio).

\subsubsection{UC-18: Elimina Spesa}

\ucblock{Utente}{Esiste almeno una spesa}{La spesa è eliminata e i totali sono aggiornati}

\noindent Flusso analogo a UC-3 (Elimina Viaggio).

% ─── UC-19 ───────────────────────────────────────────────────────────────────
\subsubsection{UC-19: Visualizza e Filtra Spese}

\ucblock{Utente}{L'utente si trova nella sezione Spese di un viaggio}{Le spese filtrate sono visualizzate con i totali aggiornati}

\noindent\textbf{Flusso normale:}
\begin{enumerate}[leftmargin=1.5em, itemsep=1pt]
  \item Il sistema mostra tutte le spese del viaggio con i totali previsto ed effettivo
  \item L'utente seleziona un filtro per categoria e/o per tipo (Prevista/Effettiva)
  \item Il sistema aggiorna la lista e ricalcola i totali in tempo reale
\end{enumerate}

% ─── UC-20 ───────────────────────────────────────────────────────────────────
\subsubsection{UC-20: Visualizza Statistiche e Riepilogo}

\ucblock{Utente}{L'utente accede alla schermata Statistiche}{Le statistiche aggregate sono visualizzate}

\noindent\textbf{Flusso normale:}
\begin{enumerate}[leftmargin=1.5em, itemsep=1pt]
  \item Il sistema calcola e visualizza il numero totale di viaggi suddivisi per stato
  \item Il sistema mostra il numero di attività pianificate e completate
  \item Il sistema mostra la spesa totale prevista e quella effettiva con il confronto
    rispetto al budget
  \item Il sistema visualizza la distribuzione delle spese per categoria tramite grafico
  \item Il sistema mostra il progresso complessivo delle checklist
  \item I dati vengono ricalcolati ad ogni accesso alla schermata
\end{enumerate}

% ─── UC-21 ───────────────────────────────────────────────────────────────────
\subsubsection{UC-21: Timeline Giornaliera (Feature Avanzata)}

\ucblock{Utente}{L'utente si trova nel dettaglio di una tappa}{Le attività della tappa sono visualizzate in una timeline oraria}

\noindent\textbf{Flusso normale:}
\begin{enumerate}[leftmargin=1.5em, itemsep=1pt]
  \item L'utente seleziona la vista ``Timeline'' dalla schermata della tappa
  \item Il sistema visualizza le attività con orario disposte verticalmente su una
    scala temporale dalle ore di inizio della prima attività alle ore di fine dell'ultima
  \item Le attività prive di orario sono mostrate in una sezione separata ``Non pianificate''
  \item Selezionando un'attività nella timeline, l'utente accede al suo dettaglio
\end{enumerate}

% ─── UC-22 ───────────────────────────────────────────────────────────────────
\subsubsection{UC-22: Duplica Tappa (Feature Avanzata)}

\ucblock{Utente}{Esiste almeno una tappa nel viaggio corrente}{Una copia della tappa con tutte le sue attività è creata}

\noindent\textbf{Flusso normale:}
\begin{enumerate}[leftmargin=1.5em, itemsep=1pt]
  \item L'utente seleziona ``Duplica tappa'' dal menu della tappa
  \item Il sistema chiede la data da assegnare alla tappa duplicata
  \item L'utente inserisce la data e conferma
  \item Il sistema crea una copia della tappa con nuovi identificatori univoci e
    la data selezionata; tutte le attività associate vengono duplicate con lo stato
    reimpostato a ``Da svolgere''
  \item La nuova tappa appare nell'itinerario nella posizione corretta per data
\end{enumerate}

\noindent\textbf{Flussi alternativi:}
\begin{itemize}[leftmargin=1.5em, itemsep=1pt]
  \item 3a. Data non inserita: il sistema segnala l'errore, l'esecuzione riprende dal passo 2.
\end{itemize}

% =============================================================================
\section{Tracciabilità}
% =============================================================================

\subsection{Matrice di Tracciabilità}

La matrice collega i requisiti ai casi d'uso che li implementano, ai componenti software
che li realizzano e agli esempi di test che ne verificano il corretto funzionamento.
La matrice viene aggiornata durante lo sviluppo man mano che i componenti vengono
definiti e i test scritti.

\begingroup
\small
\setlength{\tabcolsep}{4pt}
\begin{longtable}{@{}L{1.8cm}L{3.5cm}L{3.5cm}L{5.2cm}@{}}
  \toprule
  \textbf{Requisito} & \textbf{Casi d'Uso} & \textbf{Componenti} & \textbf{Casi di Test (esempi)} \\
  \midrule
  \endfirsthead
  \toprule
  \textbf{Requisito} & \textbf{Casi d'Uso} & \textbf{Componenti} & \textbf{Casi di Test (esempi)} \\
  \midrule
  \endhead
  \bottomrule
  \endlastfoot

  BF-1 & UC-1 -- UC-5 & TripListScreen, TripFormScreen, TripDetailScreen, TripStore &
    testCreaTripValido, testModificaTrip, testEliminaTrip, testFiltroPerStato \\
  \rowcolor{rowgray}
  BF-2 & UC-6 -- UC-9 & StageListScreen, StageFormScreen, StageDetailScreen, StageStore &
    testCreaStageValido, testEliminaStageACascata, testOrdinamentoPerData \\
  BF-3 & UC-10 -- UC-13 & ActivityListScreen, ActivityFormScreen, ActivityStore &
    testCreaActivity, testModificaStato, testFiltroPerCategoria \\
  \rowcolor{rowgray}
  BF-4 & UC-14 -- UC-15 & ChecklistScreen, ChecklistDetailScreen, ChecklistStore &
    testCreaChecklist, testToggleItem, testProgressoChecklist \\
  BF-5 & UC-16 -- UC-19 & ExpenseListScreen, ExpenseFormScreen, ExpenseStore &
    testCreaExpense, testImportoNonValido, testTotaleSpese \\
  \rowcolor{rowgray}
  BF-6 & UC-5, UC-13, UC-19, UC-20 & StatsScreen, SearchService &
    testFiltroRealTime, testCalcoloStatistiche \\
  \midrule
  IF-1.1 & UC-1 & TripFormScreen, TripStore &
    testCreaTripValido, testTitoloVuoto, testDateInvalide \\
  \rowcolor{rowgray}
  IF-1.2 & UC-2 & TripFormScreen, TripStore &
    testModificaTripSuccesso, testDataFineAntecedente \\
  IF-1.3 & UC-3 & TripDetailScreen, TripStore &
    testEliminaTrip, testEliminazioneACascata \\
  \rowcolor{rowgray}
  IF-1.4 & UC-4 & TripListScreen &
    testOrdinamentoPerData, testStatoVuoto \\
  IF-1.5 & UC-5 & TripListScreen, SearchService &
    testRicercaDestinazione, testFiltroStato, testNessunRisultato \\
  \rowcolor{rowgray}
  IF-2.1 & UC-6 & StageFormScreen, StageStore &
    testCreaStageValido, testTitoloVuoto \\
  IF-2.2 & UC-7 & StageFormScreen, StageStore &
    testModificaStage \\
  \rowcolor{rowgray}
  IF-2.3 & UC-8 & StageDetailScreen, StageStore &
    testEliminaStageConAttivita \\
  IF-2.4 & UC-9 & StageListScreen &
    testOrdinamentoStage, testNumeroAttivita \\
  \rowcolor{rowgray}
  IF-3.1 & UC-10 & ActivityFormScreen, ActivityStore &
    testCreaActivityValida, testCostoNegativo \\
  IF-3.2 & UC-11 & ActivityFormScreen, ActivityStore &
    testModificaStato \\
  \rowcolor{rowgray}
  IF-3.3 & UC-12 & ActivityListScreen, ActivityStore &
    testEliminaActivity \\
  IF-3.4 & UC-13 & ActivityListScreen &
    testFiltroPerStato, testFiltroPerCategoria \\
  \rowcolor{rowgray}
  IF-4.1 & UC-14 & ChecklistScreen, ChecklistStore &
    testCreaChecklist, testProgressoChecklist, testToggleCompletato \\
  IF-4.2 & UC-15 & TripInfoScreen &
    testSalvaNotaLibera, testModificaNota \\
  \rowcolor{rowgray}
  IF-5.1 & UC-16 & ExpenseFormScreen, ExpenseStore &
    testCreaExpense, testImportoZero \\
  IF-5.2 & UC-17, UC-18 & ExpenseFormScreen, ExpenseListScreen &
    testModificaExpense, testEliminaExpense \\
  \rowcolor{rowgray}
  IF-5.3 & UC-19 & ExpenseListScreen &
    testTotalePrevisto, testTotaleEffettivo, testFiltroCategoria \\
  IF-6.1 & UC-20 & StatsScreen &
    testCalcoloStatistiche, testDistribuzionePerCategoria \\
  \midrule
  \rowcolor{rowgray}
  DF-1.1 & UC-1, UC-2, UC-3 & Trip (tipo TypeScript), TripStore &
    testCostruttoreTrip, testValidazioneDate \\
  DF-1.2 & UC-6, UC-7, UC-8 & Stage (tipo TypeScript), StageStore &
    testCostruttoreStage \\
  \rowcolor{rowgray}
  DF-1.3 & UC-10, UC-11, UC-12 & Activity (tipo TypeScript), ActivityStore &
    testCostruttoreActivity, testStatoDefault \\
  DF-1.4 & UC-16, UC-17, UC-18 & Expense (tipo TypeScript), ExpenseStore &
    testCostruttoreExpense \\
  \rowcolor{rowgray}
  DF-1.5 & UC-14 & Checklist, ChecklistItem (tipi TypeScript), ChecklistStore &
    testToggleItem, testProgressoCalculation \\
  DF-2.1 & UC-1 -- UC-22 & StorageService &
    testSalvataggioPersistente, testCaricamentoAvvio \\
  \midrule
  \rowcolor{rowgray}
  FC-1 & UC-1 -- UC-22 & Tutti i componenti &
    testGestioneEccezione, testNoCrashInputVuoto \\
  FC-2 & UC-3, UC-8 & TripStore, StageStore &
    testEliminazioneACascata, testRiferimentoNullo \\
  \rowcolor{rowgray}
  FC-3 & UC-1, UC-6, UC-10, UC-16 & Tutti i FormScreen &
    testCampiObbligatori, testRangeValori \\
  FC-7 & UC-1 -- UC-22 & StorageService &
    testFunzionamentoOffline \\
\end{longtable}
\endgroup

\subsection{Legenda}

\begin{tabular}{@{}ll@{}}
  \toprule
  \textbf{Prefisso} & \textbf{Significato} \\
  \midrule
  BF & Business Flow \\
  IF & Funzionalità Individuale \\
  DF & Requisito sui Dati \\
  UI & Requisito di Interfaccia Utente \\
  FC & Fattore di Qualità (requisito non funzionale) \\
  UC & Caso d'Uso \\
  \bottomrule
\end{tabular}

% =============================================================================
\appendix
\section{Glossario}
% =============================================================================

\begin{tabular}{@{}p{3.5cm}p{10.5cm}@{}}
  \toprule
  \textbf{Termine} & \textbf{Definizione} \\
  \midrule
  Viaggio & Entità principale dell'app. Rappresenta un piano di viaggio con date, destinazione, budget e dati ad esso collegati. \\
  Tappa & Suddivisione di un viaggio corrispondente a una giornata o a una specifica città o area geografica. \\
  Attività & Evento o impegno pianificato durante il viaggio (visita, pasto, spostamento, ecc.), associato a un viaggio e opzionalmente a una tappa. \\
  Checklist & Lista di elementi da preparare o spuntare prima o durante il viaggio. \\
  Spesa prevista & Costo stimato prima della partenza o dell'attività. \\
  Spesa effettiva & Costo realmente sostenuto durante il viaggio. \\
  Budget & Importo massimo di spesa pianificato per il viaggio. \\
  Timeline & Vista grafica delle attività di una tappa disposta su una scala oraria verticale. \\
  Duplicazione tappa & Feature avanzata che crea una copia di una tappa (con le relative attività) su una nuova data. \\
  Storage locale & Meccanismo di persistenza dei dati sul dispositivo (AsyncStorage) che non richiede connessione di rete. \\
  CRUD & Create, Read, Update, Delete -- le quattro operazioni fondamentali di gestione dei dati. \\
  Offline-first & Principio progettuale per cui l'app funziona completamente senza connessione di rete. \\
  Store & Modulo di gestione dello stato globale implementato con Zustand. \\
  \bottomrule
\end{tabular}

\end{document}
